\documentclass[reprint, amsmath,amssymb, amsthm, aip, cp]{revtex4-2}

% LaTeX Template by Aditya Rao
% Use this for organizing images for writing into laboratory notebooks.
%=================== HEADER =====================%
\usepackage[left=0.7in, right=0.7in, top=0.7in, bottom=0.7in]{geometry}

\usepackage[T1]{fontenc}
\usepackage[utf8]{inpute nc}
\usepackage{setspace}
\usepackage{fancyhdr}
\usepackage{lipsum}


% Other Preamble
\usepackage[leftmargin=1cm, rightmargin=1cm]{mdframed}
\usepackage[colorlinks, allcolors=black]{hyperref}


% Code Formatting
\usepackage{listings}

% Figures & Drawings
\usepackage{graphicx}
\usepackage[export]{adjustbox}
\usepackage{tikz}
\usepackage{float}
\usepackage{caption}
\usetikzlibrary{patterns}
\usetikzlibrary{calc}

% Physics & Mathematics
\usepackage{physics}
\usepackage{mathtools}
%================================================%

%=================== SETUP ======================%
% Page styling
\pagestyle{fancy}
\pagenumbering{arabic}

% Listings Setup
\lstset{
	basicstyle=\footnotesize,
	numbers=left,
	stepnumber=1,
	showstringspaces=false,
	tabsize=1,
	breaklines=true,
	breakatwhitespace=false
}

% Date Formatting
\def\mydate{\leavevmode\hbox{\the\year-\twodigits\month-\twodigits\day}}
\def\twodigits#1{\ifnum#1<10 0\fi\the#1}
%================================================%

%=================== MACROS =====================%
% The main stuff
\let\origthefigure\thefigure

\newcommand{\thedate}{00.00}
\newcommand{\figbig}{0.9\linewidth}
\newcommand{\figsml}{0.7\linewidth}
\newcommand{\newdate}[1]{\newpage\renewcommand{\thedate}{#1}\section*{Date: \thedate}}
\newcommand{\updatefig}[1]{\renewcommand{\thefigure}{\textbf{\thedate #1.\origthefigure}}\setcounter{figure}{0}}

\newcommand{\updatetab}[1]{\renewcommand{\thefigure}{\textbf{\thedate #1.\origthefigure}}\setcounter{table}{0}}

\newcommand{\updatecountfig}[1]{\setcounter{figure}{#1}}
\newcommand{\updatecounttab}[1]{\setcounter{table}{#1}}

\newcommand{\mdfigure}[2]{
	\begin{mdframed}
	\begin{figure}[H]
		\centering
		\captionsetup{width=\figsml}
		\includegraphics[width=\figsml, max height=0.4\textheight]{#1}
		\caption{#2}
	\end{figure}
	\end{mdframed}
}

\newcommand{\mdfiguretex}[2]{
	\begin{mdframed}
	\begin{figure}[H]
		\centering
		\captionsetup{width=\figsml}
		\input{#1}
		\caption{#2}
	\end{figure}
	\end{mdframed}
}

\newcommand{\mdtable}[2]{
	\begin{mdframed}
	\begin{figure}[H]
		\centering
		\captionsetup{width=\figsml}
		\includegraphics[width=\figsml, max height=0.4\textheight]{#1}
		\caption{#2}
	\end{figure}
	\end{mdframed}
}


\newcommand{\mdcode}[1]{
	\begin{mdframed}
	\lstset{frame=shadowbox,escapechar=`,linewidth=6cm}
	\lstinputlisting[language=python]{#1}
	\end{mdframed}
}

% I like the Griffiths' script r
%Script r (vector)
\newcommand{\brc}{\resizebox{!}{1.3ex}{
    \begin{tikzpicture}[>=round cap]
        \clip (0.085em,-0.1ex) rectangle (0.61em,0.875ex);
        \draw [line width=.2ex, <->, rounded corners=0.13ex] (0.1em,0.1ex) .. controls (0.24em,0.4ex) .. (0.35em,0.8ex) .. controls (0.29em,0.725ex) .. (0.25em,0.6ex) .. controls (0.7em,0.8ex) and (0.08em,-0.4ex) .. (0.55em,0.25ex);
    \end{tikzpicture}
}}

%Script r with a hat (unit vector)
\newcommand{\hrc}{\hat{\brc}}

% References and Sourcing
\newcommand{\figref}[1]{Fig. \ref{#1}}
\newcommand{\source}[1]{\caption*{Source: {#1}} }
%================================================%

%==================== INFO =======================%
\newcommand{\name}{Aditya K. Rao}
\newcommand{\expname}{QIE}
%================================================%

\begin{document}
\nocite{*}
\lhead{\name\vspace{0.1cm}}
\chead{\textbf{Advanced Physics Laboratory} \\ \expname\vspace{0.1cm}}
\rhead{\textbf{LAST UPDATED:} \mydate\vspace{0.1cm}}


\newdate{00.00}
\updatefig{X}
\mdfigure{example-image-a}{Example image for templating.}
\mdfigure{example-image-b}{Long caption. \lipsum[1].}


\newdate{03.11}
\updatefig{A}
\mdfigure{./images/03.11/expsetup1.jpg}{Intiial experimental setup prior to any changes. All components seem to be in place. Noticed no caps on extra APD detectors.}


\mdfigure{./images/03.11/expsetup2.jpg}{Experimental setup within black box. Laser and powersupply and all other relavent components outlined in P03.11A turned off prior to opening.}


\mdfigure{./images/03.11/tauval.png}{Initial fit of coincidence values and expected statistical correction. Slope found is given by $\tau = 0.0243 \pm 0.0006$ with a $\chi^2_{\text{red}} = 0.03$.}

\newdate{03.14}
\updatefig{A}

\mdfigure{../Images/03.14/coincwithres.png}{Same fit of data as in Fig. 03.11A.3 but with residuals.}

\newdate{03.18}
\updatefig{A}

\mdfigure{../Images/03.18/CcalcRes.png}{Plot of series of $N(90^\circ, 0^\circ)$ and $N(0^\circ, 90^\circ)$ data points to illustrate variation in data.}

\newdate{03.19}
\updatefig{D}
\updatecountfig{1}
\mdfigure{../Images/03.19/spdc.png}{Diagram of Spontaneous Paramtric Down Conversion (SPDC) process. \textbf{Source:} Applications of single-photon technology - Scientific Figure on ResearchGate. Available from: \href{https://www.researchgate.net/figure/The-SPDC-cones-a-Type-II-process-b-Polarization-entangled-photon-pair_fig10_360782305}{https://www.researchgate.net/figure/The-SPDC-cones-a-Type-II-process-b-Polarization-entangled-photon-pair\_fig10\_360782305} [accessed 19 Mar 2025].}

\updatefig{E}
\updatecountfig{2}
\mdfiguretex{../Images/03.19/expsetup.tex}{Experimental setup ($405\,\text{nm}$).}

\newdate{03.26}
\updatefig{A}
\mdfigure{../Images/UCat/HWP.png}{Half Waveplate and Quartz Plate apparatus.}

\updatefig{B}
\mdfigure{../Images/UCat/LASER.jpg}{$405\,\text{nm}$ Diode Laser (beginning of setup.)}

\updatefig{C}
\mdfigure{../Images/UCat/APD.png}{Avalanche Photodiode (APD) detector and Polarizer Setup.}

\newdate{03.22}
\updatefig{A}
\mdfigure{../Images/03.22/BBO.png}{The extraordinary and ordinary refractive indices (and their difference) as a function of direction for beta barium borate (BBO) at $\lambda=500\,\text{nm}$ ($n_o=1.68$ and $n_e = 1.56$) \cite{pwoptics}.}


\newdate{03.27}
\updatefig{A}
\mdfigure{../Images/04.02/NA90B112p5_EA0B22p5.png}{}
\mdfigure{../Images/04.02/NA0B22p5_EA0B22p5.png}{}
\mdfigure{../Images/04.02/NA90B22p5_EA0B22p5.png}{}
\mdfigure{../Images/04.02/NA0B112p5_EA0B22p5.png}{}


\updatefig{B}
\mdfigure{../Images/04.02/NA45B112p5_EAm45B22p5.png}{}
\mdfigure{../Images/04.02/NAm45B22p5_EAm45B22p5.png}{}
\mdfigure{../Images/04.02/NA45B22p5_EAm45B22p5.png}{}
\mdfigure{../Images/04.02/NAm45B112p5_EAm45B22p5.png}{}

\updatefig{C}
\mdfigure{../Images/04.02/NA90B67p5_EA0Bm22p5.png}{}
\mdfigure{../Images/04.02/NA0Bm22p5_EA0Bm22p5.png}{}
\mdfigure{../Images/04.02/NA90Bm22p5_EA0Bm22p5.png}{}
\mdfigure{../Images/04.02/NA0B67p5_EA0Bm22p5.png}{}

\updatefig{E}
\mdfigure{../Images/04.02/NA45B67p5_EAm45Bm22p5.png}{}
\mdfigure{../Images/04.02/NAm45Bm22p5_EAm45Bm22p5.png}{}

\newdate{03.25}
\updatefig{A}
\mdfigure{../Images/04.02/NA45Bm22p5_EAm45Bm22p5.png}{}
\mdfigure{../Images/04.02/NAm45B67p5_EAm45Bm22p5.png}{}


\newdate{03.30}
\updatefig{B}
\mdfigure{../Images/03.30/aspect.png}{Illustration between delay time and coincidence rate \cite{aspect_2025} known as the \textit{Time Delay Spectrum}. The flat background corresponds with accidental coincidences between uncorrelated photons emitted by different atoms and scales with the square, $N^2$, of the cascade rate. The peak, whos area scales with $N$, corresponds with the correlated photons and gives the coincidences to be measured.}


\newdate{04.03}
\updatefig{A}
\mdfigure{../Images/04.02/sipe.png}{Space like seperated regions as $S_A$ is not in the backward lightcone of $S_B$. Figure from Prof. Sipe's PHY491 notes \cite{sipe}.}

\updatefig{B}
\mdfigure{../Images/04.02/RAO1.png}{(a) Two events/interactions  at the same spacetime location are local. (b) Two time-like events separated by a distance preserve Einstein's locality condition. (c) Two space-like seperated events are nonlocal. Figure and caption from \cite{rao2024investigating}.}

\updatefig{C}
\mdfigure{../Images/04.02/RAO2.png}{Quantum teleportation circuit as described (with modification) by Nielson and Chuang \cite{nielsen_chuang_2010}. The first Hadamard ($\hat{H}$) and CNOT gates describe the creation of a Bell state. Figure and caption from \cite{rao2024investigating}.}

\newdate{04.04}
\updatefig{A}
\mdfigure{../Images/UCat/APD.png}{APD and polarizer setup. This is the same figure as Fig 03.2C.1. The main thing to note is that the fibers are not adequately spaced before the measurement devices.}


\clearpage
\appendix
\bibliography{References.bib}
\end{document}